\chapter{Результаты}
\label{chap:results}

В таблице~\ref{tab:val_metrics} приведены валидационные метрики шести моделей относительно порогов приёмки. После регрессионной стадии порог по PSNR проходят два глубоких генератора, а по SSIM не проходит ни один. После состязательного дообучения все три модели укладываются во все пороги, причём PSNR и SSIM растут одновременно. Наилучшие значения по всем метрикам показывает RFDN GAN, остающийся при этом самым компактным: около 3 МБ в формате ONNX против 5.8 МБ у EDSR.

\begin{table}[h]
\captionsetup{justification=centering}
\centering
\caption{Валидационные метрики относительно порогов приёмки.}
\label{tab:val_metrics}
\small
\setlength{\tabcolsep}{6pt}
\begin{tabular}{l cccc}
\toprule
Конфигурация
  & PSNR$\uparrow$ & SSIM$\uparrow$ & GMAE$\downarrow$ & ERGAS$\downarrow$ \\
Порог
  & $\geq 45$ & $\geq 0.55$ & $\leq 0.002$ & $\leq 0.4$ \\
\midrule
\multicolumn{5}{l}{Регрессионная стадия} \\
\midrule
MSRResNet & 44.92          & \textbf{0.559} & \textbf{0.00195} & \textbf{0.380} \\
EDSR      & \textbf{46.92} & 0.548          & \textbf{0.00064} & \textbf{0.272} \\
RFDN      & \textbf{46.11} & 0.505          & \textbf{0.00083} & \textbf{0.310} \\
\midrule
\multicolumn{5}{l}{Состязательное дообучение} \\
\midrule
MSRResNet GAN & \textbf{45.99} & \textbf{0.590} & \textbf{0.00165} & \textbf{0.353} \\
EDSR GAN      & \textbf{47.71} & \textbf{0.606} & \textbf{0.00062} & \textbf{0.257} \\
RFDN GAN      & \textbf{47.77} & \textbf{0.634} & \textbf{0.00070} & \textbf{0.264} \\
\bottomrule
\end{tabular}
\end{table}